% GENERATED FILE. Edit canonical lessons, then run npm run generate:books.
% canonical-source-hash: fnv1a64:811b2ae812ee6b00
% canonical-lessons: TE-C428-nindincu, TE-C428-avamanincu, TE-C428-vismarincu, TE-C428-lopam, TE-C428-savarincu

\chapter{Blame and Correction}
\label{ch:te-blame-and-correction}
\hlchaptermodality{428}

\begin{chapteropening}
\textbf{By the end of this chapter:} \emph{I can blame, insult, overlook and put things right.}

Complete the last lesson of chapter 428: I can blame, insult, overlook and put things right.
\end{chapteropening}

\section[\texorpdfstring{\te{నిందించు} (nindiñcu)}{నిందించు (nindiñcu)}]{\te{నిందించు} (nindiñcu) — to reproach; to blame}

\label{lesson:TE-C428-nindincu}



\begin{quote}
Before the new one: say the Telugu for to satisfy, then the Telugu for a mediator.
\end{quote}

\subsection*{You'll want to know: \te{నిందించు}}
\textbf{\te{నిందించు}} — \emph{nindiñcu} — \textquotedblleft{}to reproach; to blame\textquotedblright{}.

\begin{quote}
\te{తన} \te{తప్పుకు} \te{అతను} \te{ఇతరులను} \te{నిందించాడు}.
\end{quote}

\begin{quote}
\emph{tana tappuku atanu itarulanu nindiñcāḍu.}
\end{quote}

\begin{quote}
He blamed others for his own mistake.
\end{quote}

Past (I …): \te{నిందించాను}, \emph{nindiñcānu}.

The verb of \te{నింద}, blame, which you know.

\begin{cousinweb}
\textbf{\te{నిందించు}} (\emph{nindiñcu}) is built on Sanskrit \emph{nindā}, \textquotedblleft{}blame\textquotedblright{} (CDIAL 7212), with the Telugu verb ending -\te{ఇంచు}; \te{నింద} is the same word.
\end{cousinweb}

\begin{grammarlens}[title={What you've built}]
One more word for saying what happens if.
\end{grammarlens}

\subsection*{Your turn}
\begin{itemize}
\raggedright
  \item \emph{Say it:} \emph{nindiñcu}
  \item \emph{Say it:} \emph{nindiñcu}, once more
  \item \emph{Say it:} \emph{tana tappuku atanu itarulanu nindiñcāḍu.} — He blamed others for his own mistake.
  \item \emph{Recall:} say the Telugu for to satisfy, then the Telugu for a mediator, then say \emph{nindiñcu} again
\end{itemize}

\begin{tcolorbox}[breakable,colback=teal!4,colframe=teal!35!black,title={Before you move on}]
What does \te{నిందించు} mean? (\textquotedblleft{}To reproach; to blame\textquotedblright{}.) Say it once more.
\end{tcolorbox}
\section[\texorpdfstring{\te{అవమానించు} (avamāniñcu)}{అవమానించు (avamāniñcu)}]{\te{అవమానించు} (avamāniñcu) — to insult; to humiliate}

\label{lesson:TE-C428-avamanincu}



\begin{quote}
Before the new one: say the Telugu for a mediator, then the Telugu for to reproach.
\end{quote}

\subsection*{You'll want to know: \te{అవమానించు}}
\textbf{\te{అవమానించు}} — \emph{avamāniñcu} — \textquotedblleft{}to insult; to humiliate\textquotedblright{}.

\begin{quote}
\te{అందరి} \te{ముందు} \te{అతను} \te{నన్ను} \te{అవమానించాడు}.
\end{quote}

\begin{quote}
\emph{andari mundu atanu nannu avamāniñcāḍu.}
\end{quote}

\begin{quote}
He insulted me in front of everyone.
\end{quote}

Past (I …): \te{అవమానించాను}, \emph{avamāniñcānu}.

The person insulted takes -\te{ని}: \te{నన్ను} \te{అవమానించాడు}.

\begin{cousinweb}
\textbf{\te{అవమానించు}} (\emph{avamāniñcu}) is built on Sanskrit \emph{avamāna}, \textquotedblleft{}disrespect, contempt, dishonour\textquotedblright{} (Monier-Williams), with the Telugu verb ending -\te{ఇంచు}.
\end{cousinweb}

\begin{grammarlens}[title={What you've built}]
One more word for saying what happens if.
\end{grammarlens}

\subsection*{Your turn}
\begin{itemize}
\raggedright
  \item \emph{Say it:} \emph{avamāniñcu}
  \item \emph{Say it:} \emph{avamāniñcu}, once more
  \item \emph{Say it:} \emph{andari mundu atanu nannu avamāniñcāḍu.} — He insulted me in front of everyone.
  \item \emph{Recall:} say the Telugu for a mediator, then the Telugu for to reproach, then say \emph{avamāniñcu} again
\end{itemize}

\begin{tcolorbox}[breakable,colback=teal!4,colframe=teal!35!black,title={Before you move on}]
What does \te{అవమానించు} mean? (\textquotedblleft{}To insult; to humiliate\textquotedblright{}.) Say it once more.
\end{tcolorbox}
\section[\texorpdfstring{\te{విస్మరించు} (vismariñcu)}{విస్మరించు (vismariñcu)}]{\te{విస్మరించు} (vismariñcu) — to ignore; to overlook}

\label{lesson:TE-C428-vismarincu}



\begin{quote}
Before the new one: say the Telugu for to reproach, then the Telugu for to insult.
\end{quote}

\subsection*{You'll want to know: \te{విస్మరించు}}
\textbf{\te{విస్మరించు}} — \emph{vismariñcu} — \textquotedblleft{}to ignore; to overlook\textquotedblright{}.

\begin{quote}
\te{ప్రభుత్వం} \te{మా} \te{ఫిర్యాదును} \te{విస్మరించింది}.
\end{quote}

\begin{quote}
\emph{prabhutvaṁ mā phiryādunu vismariñcindi.}
\end{quote}

\begin{quote}
The government ignored our complaint.
\end{quote}

Past (I …): \te{విస్మరించాను}, \emph{vismariñcānu}.

Formal; in talk, \te{పట్టించుకోలేదు}, took no notice.

\begin{cousinweb}
\textbf{\te{విస్మరించు}} (\emph{vismariñcu}) is built on Sanskrit \emph{vismaraṇa}, \textquotedblleft{}the act of forgetting, oblivion\textquotedblright{} (Monier-Williams), with the Telugu verb ending -\te{ఇంచు}.
\end{cousinweb}

\begin{grammarlens}[title={What you've built}]
One more word for saying what happens if.
\end{grammarlens}

\subsection*{Your turn}
\begin{itemize}
\raggedright
  \item \emph{Say it:} \emph{vismariñcu}
  \item \emph{Say it:} \emph{vismariñcu}, once more
  \item \emph{Say it:} \emph{prabhutvaṁ mā phiryādunu vismariñcindi.} — The government ignored our complaint.
  \item \emph{Recall:} say the Telugu for to reproach, then the Telugu for to insult, then say \emph{vismariñcu} again
\end{itemize}

\begin{tcolorbox}[breakable,colback=teal!4,colframe=teal!35!black,title={Before you move on}]
What does \te{విస్మరించు} mean? (\textquotedblleft{}To ignore; to overlook\textquotedblright{}.) Say it once more.
\end{tcolorbox}
\section[\texorpdfstring{\te{లోపం} (lōpaṁ)}{లోపం (lōpaṁ)}]{\te{లోపం} (lōpaṁ) — a flaw; a shortcoming}

\label{lesson:TE-C428-lopam}



\begin{quote}
Before the new one: say the Telugu for to insult, then the Telugu for to ignore.
\end{quote}

\subsection*{You'll want to know: \te{లోపం}}
\textbf{\te{లోపం}} — \emph{lōpaṁ} — \textquotedblleft{}a flaw; a shortcoming\textquotedblright{}.

\begin{quote}
\te{ఆ} \te{ప్రణాళికలో} \te{ఒక} \te{పెద్ద} \te{లోపం} \te{ఉంది}.
\end{quote}

\begin{quote}
\emph{ā praṇāḷikalō oka pedda lōpaṁ undi.}
\end{quote}

\begin{quote}
There is a big flaw in that plan.
\end{quote}

Plural: \te{లోపాలు}, \emph{lōpālu}.

Plural \te{లోపాలు}.

\begin{cousinweb}
\textbf{\te{లోపం}} (\emph{lōpaṁ}) is borrowed from Sanskrit \emph{lopa}, \textquotedblleft{}breaking, injury; neglect, violation\textquotedblright{} (Monier-Williams); Telugu uses it for a fault or a lack.
\end{cousinweb}

\begin{grammarlens}[title={What you've built}]
One more word for saying what happens if.
\end{grammarlens}

\subsection*{Your turn}
\begin{itemize}
\raggedright
  \item \emph{Say it:} \emph{lōpaṁ}
  \item \emph{Say it:} \emph{lōpaṁ}, once more
  \item \emph{Say it:} \emph{ā praṇāḷikalō oka pedda lōpaṁ undi.} — There is a big flaw in that plan.
  \item \emph{Recall:} say the Telugu for to insult, then the Telugu for to ignore, then say \emph{lōpaṁ} again
\end{itemize}

\begin{tcolorbox}[breakable,colback=teal!4,colframe=teal!35!black,title={Before you move on}]
What does \te{లోపం} mean? (\textquotedblleft{}A flaw; a shortcoming\textquotedblright{}.) Say it once more.
\end{tcolorbox}
\section[\texorpdfstring{\te{సవరించు} (savariñcu)}{సవరించు (savariñcu)}]{\te{సవరించు} (savariñcu) — to amend; to correct}

\label{lesson:TE-C428-savarincu}



\begin{quote}
Before the new one: say the Telugu for to ignore, then the Telugu for a flaw.
\end{quote}

\subsection*{You'll want to know: \te{సవరించు}}
\textbf{\te{సవరించు}} — \emph{savariñcu} — \textquotedblleft{}to amend; to correct\textquotedblright{}.

\begin{quote}
\te{ఉపాధ్యాయుడు} \te{నా} \te{తప్పులను} \te{సవరించాడు}.
\end{quote}

\begin{quote}
\emph{upādhyāyuḍu nā tappulanu savariñcāḍu.}
\end{quote}

\begin{quote}
The teacher corrected my mistakes.
\end{quote}

Past (I …): \te{సవరించాను}, \emph{savariñcānu}.

The noun is \te{సవరణ}, an amendment, a correction.

\begin{cousinweb}
\textbf{\te{సవరించు}} (\emph{savariñcu}) is native Dravidian: DEDR 2342 lists Telugu \emph{savarincu}, \textquotedblleft{}to adjust, to trim, to put in order\textquotedblright{}, and \emph{savaraṇa}, \textquotedblleft{}correction, setting right\textquotedblright{}.
\end{cousinweb}

\begin{grammarlens}[title={What you've built}]
One more word for saying what happens if.
\end{grammarlens}

\subsection*{Your turn}
\begin{itemize}
\raggedright
  \item \emph{Say it:} \emph{savariñcu}
  \item \emph{Say it:} \emph{savariñcu}, once more
  \item \emph{Say it:} \emph{upādhyāyuḍu nā tappulanu savariñcāḍu.} — The teacher corrected my mistakes.
  \item \emph{Recall:} say the Telugu for to ignore, then the Telugu for a flaw, then say \emph{savariñcu} again
\end{itemize}

\begin{tcolorbox}[breakable,colback=teal!4,colframe=teal!35!black,title={Before you move on}]
What does \te{సవరించు} mean? (\textquotedblleft{}To amend; to correct\textquotedblright{}.) Say it once more.
\end{tcolorbox}
